\documentclass[11pt]{article}
\usepackage[margin=1in]{geometry}
\usepackage{amsmath,amssymb,amsthm}
\usepackage{hyperref}
\usepackage{xurl}

\newtheorem{theorem}{Theorem}
\newtheorem{lemma}[theorem]{Lemma}
\newtheorem{corollary}[theorem]{Corollary}
\theoremstyle{definition}
\newtheorem{definition}[theorem]{Definition}
\theoremstyle{remark}
\newtheorem{remark}[theorem]{Remark}

\let\oldus\_
\renewcommand{\_}{\oldus\allowbreak}

\title{A disproof of Conjecture 00000005640\\[2pt]
\large An intersection-preserving redrawing cannot change the intersection graph}
\author{Jackmeson1}
\date{}

\begin{document}
\sloppy
\maketitle

\section{The conjecture}

Conjecture 00000005640 reads (English version; the Chinese version says the same, with
``the separation is realized by an explicit redrawing pair with intersections preserved'' for the
last clause):

\begin{quote}
\emph{Definition:} The intersection spectrum and the intersection graph are two layers.
\emph{Conjecture:} There exist two convex-set families with the same intersection spectrum but
non-isomorphic intersection graphs, and the separation is realized by an explicit redrawing pair
preserving intersections.
\end{quote}

The claim is a conjunction: two convex families with (A) the same intersection spectrum and (B)
non-isomorphic intersection graphs, where (C) this separation is realized by an explicit
\emph{redrawing pair preserving intersections}. We show that (C) and (B) are incompatible, whatever
(A) means. Hence the conjecture is false.

\section{Reading and definitions}

\begin{definition}[Intersection graph]
For an indexed family $F=(F_i)_{i\in I}$ of sets, the intersection graph $\Gamma(F)$ has vertex set
$I$, and distinct $i\neq j$ are adjacent iff $F_i\cap F_j\neq\emptyset$. This is the standard
notion (see e.g.\ \url{https://en.wikipedia.org/wiki/Intersection_graph}).
\end{definition}

\begin{definition}[Redrawing pair preserving intersections]
Two families $F=(F_i)_{i\in I}$ and $G=(G_k)_{k\in K}$ form a \emph{redrawing pair preserving
intersections} if there is a bijection $\sigma\colon I\to K$ (each member $F_i$ is redrawn as
$G_{\sigma(i)}$) such that for all $i\neq j$:
\[ F_i\cap F_j\neq\emptyset \iff G_{\sigma(i)}\cap G_{\sigma(j)}\neq\emptyset. \]
\end{definition}

A convex family in the conjecture is a family of convex subsets of $\mathbb{R}^d$; the two families
may live in different dimensions $d,d'$, and the index sets may be finite or infinite. The term
``intersection spectrum'' is not defined in the source. We leave it as an \emph{arbitrary}
relation between families, so the disproof does not depend on its meaning.

We also treat two other readings of ``preserving intersections'':
\begin{itemize}
\item \emph{Nerve reading:} for every finite $S\subseteq I$, $\bigcap_{i\in S}F_i\neq\emptyset$ iff
$\bigcap_{i\in S}G_{\sigma(i)}\neq\emptyset$. Taking $S=\{i,j\}$ gives the pairwise reading.
\item \emph{One-sided reading:} for $i\neq j$, $F_i\cap F_j\neq\emptyset$ implies
$G_{\sigma(i)}\cap G_{\sigma(j)}\neq\emptyset$ (new intersections may appear). We treat it for finite
families whose intersection graphs have the same number of edges, i.e.\ when the ``intersection
spectrum'' determines the number of intersecting pairs.
\end{itemize}

\section{The disproof}

\begin{lemma}\label{lem:iso}
If $(F,G,\sigma)$ is a redrawing pair preserving intersections, then $\sigma$ is a graph isomorphism
$\Gamma(F)\cong\Gamma(G)$.
\end{lemma}
\begin{proof}
$\sigma$ is a bijection on vertices. For $i,j\in I$: $i\neq j$ iff $\sigma(i)\neq\sigma(j)$ since
$\sigma$ is injective, and for $i\neq j$ the intersection condition is preserved by hypothesis. So
$i\sim j$ in $\Gamma(F)$ iff $\sigma(i)\sim\sigma(j)$ in $\Gamma(G)$.
\end{proof}

\begin{theorem}\label{thm:main}
For every relation ``same intersection spectrum'', there are no convex families $F$ in
$\mathbb{R}^d$ and $G$ in $\mathbb{R}^{d'}$ that have the same intersection spectrum, have
non-isomorphic intersection graphs, and form a redrawing pair preserving intersections. Also, no
redrawing pair preserving intersections has non-isomorphic intersection graphs, so the conjecture
fails even if the separating pair and the redrawing pair are allowed to be different pairs.
\end{theorem}
\begin{proof}
By Lemma~\ref{lem:iso} the graphs of any redrawing pair preserving intersections are isomorphic,
which contradicts (B).
\end{proof}

\begin{corollary}[nerve reading]
A redrawing that preserves the nerve also gives isomorphic intersection graphs.
\end{corollary}

\begin{theorem}[one-sided reading]
Let $I,K$ be finite and let $\sigma\colon I\to K$ be a bijection that keeps every intersection of
distinct members. If $\Gamma(F)$ and $\Gamma(G)$ have the same number of edges, then $\sigma$
preserves intersections, so $\Gamma(F)\cong\Gamma(G)$.
\end{theorem}
\begin{proof}
Let $H$ be the pull-back of $\Gamma(G)$ along $\sigma$. Then $H \cong \Gamma(G)$, and
$\Gamma(F)\le H$ by hypothesis. Both have the same finite number of edges, so $\Gamma(F)=H$. Hence
$G_{\sigma(i)}\cap G_{\sigma(j)}\neq\emptyset$ implies $F_i\cap F_j\neq\emptyset$ for $i\neq j$.
\end{proof}

\begin{corollary}[one-sided reading, adjacency spectrum]\label{cor:adj}
Suppose ``same intersection spectrum'' means that the adjacency matrices $A_F$, $A_G$ of the
intersection graphs have the same characteristic polynomial. Then a one-sided redrawing of finite
families gives isomorphic intersection graphs.
\end{corollary}
\begin{proof}
$A_F$ is real symmetric. By the spectral theorem $A_F=UDU^{*}$ with $U$ unitary and $D$ the
diagonal of eigenvalues, so $\operatorname{tr}(A_F^2)=\operatorname{tr}(D^2)=\sum\lambda_i^2$. Here
the $\lambda_i$ are the roots of the characteristic polynomial, with multiplicity. Also
$\operatorname{tr}(A_F^2)=\sum_v\deg v=2|E(\Gamma(F))|$. So equal characteristic polynomials give
equal edge counts, and the previous theorem applies.
\end{proof}

\begin{remark}[non-vacuity]
The objects exist. For instance, two copies of $\mathbb{R}^1$ form a convex family, and the identity
redrawing of it preserves intersections; its intersection graph is an edge.
\end{remark}

\begin{remark}[scope]
The refuted readings are exactly these: the pairwise reading and the nerve reading, for any
meaning of ``spectrum''; and the one-sided reading for finite families whenever the spectrum
determines the number of intersecting pairs, in particular for the adjacency spectrum
(Corollary~\ref{cor:adj}). \emph{Not} refuted: the one-sided reading combined with a ``spectrum''
that does not determine the number of intersecting pairs, or with infinite families. In that case
the final clause on its own can hold (for example, two disjoint members can be moved until they
meet). The disproof uses only the bijection $\sigma$ and never uses
convexity, so it holds for arbitrary set families.
\end{remark}

\section{Lean formalization}

File \texttt{lean/Conjecture5640/Basic.lean} (Lean 4.33.1, Mathlib \texttt{v4.33.1}). The main
definitions are \texttt{interGraph} (the intersection graph, a Mathlib \texttt{SimpleGraph}),
\texttt{PreservesInter}, \texttt{PreservesNerve}, \texttt{KeepsInter}, \texttt{ConvexFamily} (an
index type together with convex subsets of \texttt{EuclideanSpace} $\mathbb{R}$ \texttt{(Fin d)}),
\texttt{IsRedrawingPair}, and \texttt{Statement} / \texttt{StatementSeparate}, both parametrized
by an arbitrary \texttt{SameSpectrum} relation. The theorems are:
\begin{itemize}
\item \texttt{isoOfPreserves}: Lemma~\ref{lem:iso}, as an explicit isomorphism
$\texttt{interGraph F}\simeq_g\texttt{interGraph G}$.
\item \texttt{redrawing\_iso}, \texttt{conjecture5640\_false},
\texttt{conjecture5640\_separate\_false}: Theorem~\ref{thm:main}.
\item \texttt{nerve\_redrawing\_iso}, \texttt{oneSided\_redrawing\_iso}: the other two readings.
\item \texttt{trace\_mul\_self\_eq\_roots}, \texttt{trace\_adj\_sq},
\texttt{card\_edges\_eq\_of\_charpoly}, \texttt{oneSided\_adjSpectrum\_iso}:
Corollary~\ref{cor:adj}.
\item \texttt{sample\_redrawing}, \texttt{sample\_adj}: non-vacuity.
\end{itemize}
\texttt{lake build} succeeds. \texttt{\#print axioms} reports only \texttt{[propext,
Classical.choice, Quot.sound]} for each of \texttt{conjecture5640\_false},
\texttt{conjecture5640\_separate\_false}, \texttt{redrawing\_iso}, \texttt{nerve\_redrawing\_iso},
\texttt{oneSided\_redrawing\_iso} and \texttt{oneSided\_adjSpectrum\_iso}. The code contains no
\texttt{sorry} and no
\texttt{native\_decide}.

The transport idea follows the accepted disproof of the sibling conjecture 00000005400 (The Last
Math Competition, GPL-3.0).

\end{document}
